%% General definitions
\documentclass{article} %% Determines the general format.
\usepackage{a4wide} %% paper size: A4.
\usepackage[utf8]{inputenc} %% This file is written in UTF-8.
%% Some editors on Windows cannot save files in UTF-8.
%% If there is a problem with special characters not showing up
%% correctly, try switching "utf8" to "latin1" (ISO 8859-1).
\usepackage[T1]{fontenc} %% Format of the resulting PDF file.
\usepackage{fancyhdr} %% Package to create a header on each page.
\usepackage{lastpage} %% Used for "Page X of Y" in the header.
                      %% For this to work, you have to call pdflatex twice.
\usepackage{enumerate} %% Used to change the style of enumerations (see below).

\usepackage{amssymb} %% Definitions for math symbols.
\usepackage{amsmath} %% Definitions for math symbols.

\usepackage{tikz}  %% Pagacke to create graphics (graphs, automata, etc.)
\usetikzlibrary{shapes,arrows,calc,positioning}
\usetikzlibrary{graphs}
\usetikzlibrary{automata} %% Tikz library to draw automata
\usetikzlibrary{arrows}   %% Tikz library for nicer arrow heads
\usetikzlibrary{positioning}    %% Tikz library for more positioning options

\usepackage{hyperref} %% hyperlinks

\usepackage{chessfss} % contained in debian package texlive-games

\usepackage{listings}   % include & display python code

\tikzset{WhiteSquare/.style = {shape=rectangle, minimum width=28, minimum
    height=28, anchor=south west}}
\tikzset{BlackSquare/.style = {shape=rectangle, fill=lightgray, minimum
    width=28, minimum height=28, anchor=south west}}

%% Left side of header
\lhead{\course\\\semester\\Exercise \homeworkNumber}
%% Right side of header
\rhead{\authorname\\Page \thepage\ of \pageref{LastPage}}
%% Height of header
\usepackage[headheight=36pt]{geometry}
%% Page style that uses the header
\pagestyle{fancy}

\newcommand{\authorname}{Benjamin Regez\\Luis Tritschler}
\newcommand{\semester}{Spring Semester 2026}
\newcommand{\course}{Foundations of Artificial Intelligence}
\newcommand{\homeworkNumber}{7}

\renewcommand{\thesection}{Exercise \homeworkNumber.\arabic{section}}


\begin{document}

\section{Arc and path consistency}
\begin{enumerate}[(a)]
    \item We apply \texttt{AC-3} on $\mathcal{C}$:
    \begin{enumerate}[1.]
        \item Iteration: \texttt{revise($\mathcal{C}, v_1, v_3$)} $\rightarrow$ 
        dom($v_3$) = $\{ g, b \}$
        \item Iteration: \texttt{revise($\mathcal{C}, v_4, v_5$)} $\rightarrow$ 
        dom($v_5$) = $\{ r, g \}$
        \item Iteration: \texttt{revise($\mathcal{C}, v_1, v_2$)} $\rightarrow$ 
        dom($v_2$) = $\{ g, b \}$
        \item Iteration: \texttt{revise($\mathcal{C}, v_4, v_3$)} $\rightarrow$ 
        dom($v_3$) = $\{ g \}$
        \item Iteration: \texttt{revise($\mathcal{C}, v_3, v_2$)} $\rightarrow$ 
        dom($v_2$) = $\{ b \}$
    \end{enumerate}
    Obtained constraint network $\mathcal{C}' =
    \langle V', \text{dom}', R'_{uv} \rangle$ with
    \begin{itemize}
        \item $V' = V$
        \item dom$'(v_1) = \{ r \}$
        \item dom$'(v_2) = \{ b \}$
        \item dom$'(v_3) = \{ g \}$
        \item dom$'(v_4) = \{ b \}$
        \item dom$'(v_5) = \{ r, g \}$
        \item $R'_{uv} = R_{uv}$
    \end{itemize}
    \item We apply \texttt{PC-2} on $\mathcal{C}$:
    \begin{enumerate}[1.]
        \item Iteration: \texttt{revise-3($\mathcal{C}, v_3, v_2, v_1$)}
        $\rightarrow R_{v_3 v_2} = \{
            \langle b, g \rangle,
            \langle g, b \rangle
        \}$
        \item Iteration: \texttt{revise-3($\mathcal{C}, v_2, v_3, v_1$)}
        $\rightarrow R_{v_2 v_3} = \{
            \langle g, b \rangle,
            \langle b, g \rangle
        \}$
        \item Iteration: \texttt{revise-3($\mathcal{C}, v_1, v_3, v_2$)}
        $\rightarrow R_{v_1 v_3} = \{
            \langle r, g \rangle,
            \langle r, b \rangle
        \}$
        \item Iteration: \texttt{revise-3($\mathcal{C}, v_1, v_2, v_3$)}
        $\rightarrow R_{v_1 v_2} = \{
            \langle r, b \rangle,
            \langle r, g \rangle
        \}$
        \item Iteration: \texttt{revise-3($\mathcal{C}, v_2, v_1, v_3$)}
        $\rightarrow R_{v_2 v_1} = \{
            \langle g, r \rangle,
            \langle b, r \rangle
        \}$
        \item Iteration: \texttt{revise-3($\mathcal{C}, v_3, v_1, v_2$)}
        $\rightarrow R_{v_3 v_1} = \{
            \langle g, r \rangle,
            \langle b, r \rangle
        \}$
    \end{enumerate}
    Obtained constraint network $\mathcal{C}''
    = \langle V'', \text{dom}'', R_{uv}'' \rangle $ with
    \begin{itemize}
        \item $V'' = V$
        \item dom$''$ = dom (with reduced domains
        $\text{dom}(v_1) = \{ r \}$ and
        $\text{dom}(v_4) = \{ b \}$)
        \item $R_{v_1, v_2}'' = \{
            \langle r, b \rangle,
            \langle r, g \rangle
        \}$
        \item $R_{v_1, v_3}'' = \{
            \langle r, g \rangle,
            \langle r, b \rangle
        \}$
        \item $R_{v_2, v_3}'' = \{
            \langle g, b \rangle,
            \langle b, g \rangle
        \}$
        \item $R_{v_2, v_5}'' = R_{v_2, v_5}$
        \item $R_{v_3, v_4}'' = R_{v_3, v_4}$
        \item $R_{v_4, v_5}'' = R_{v_4, v_5}$
    \end{itemize}
    As suggested, we implemented 7.1(b) as a python program
    instead of solving it with pen and paper.
\end{enumerate}

\newpage
\section{}

\section*{Statement on scientific integrity}



\end{document}
